\documentclass[12 pt]{article}

% Fonts/Symbols/Colors
\usepackage{fullpage}
\usepackage{amsthm}
\usepackage{amsmath}
\usepackage{xcolor}
\usepackage{amssymb}
\usepackage{verbatim}
\usepackage[utf8x]{inputenc}
\usepackage{helvet}
\usepackage{hyperref}
\hypersetup{
    colorlinks=true,
    linkcolor=blue,
    filecolor=magenta,      
    urlcolor=cyan,
    pdftitle={Overleaf Example},
    pdfpagemode=FullScreen,
    }
\renewcommand{\familydefault}{\sfdefault}

% Question counter
\newcounter{qnum}

% Question Formatting
\newcommand{\new}[1]{\bigskip \noindent \color{blue} \textbf{#1}\medskip}
\newcommand{\response}[1]{\color{black} #1 \medskip}
\newcommand{\q}[1]{\response{\emph{#1}}}
\newcommand{\Q}[1]{\stepcounter{qnum}\new{Goal {\arabic{qnum}}:} \q{#1}}

% INSERT WEEK NUMBER HERE
\newcommand{\setnum}{1}

\begin{document}

% Title Section
\begin{center}
    \Large \color{blue} CS134: Robotics Systems\color{black}
    \begin{center}\normalsize Team GLaDOS Weekly Report \setnum{}\end{center}
    \noindent{\color{black} \rule{\linewidth}{0.5mm} }
\end{center}

% Body
\noindent
\textbf{Progress}:
\begin{enumerate}
    \item The robot hardware was tested using the provided \emph{demo134} code and 
    all worked.
    \item The robot was fully assembled and is now functional.
    \item Each team member set up ssh and remote desktop on their respective machines.
    \item A team Github repository was created and setup with all team members able 
    to work on code together.
    \item Code was written to allow the robot to move to a defined starting position 
    from any initial position via a spline, and then execute a periodic waving motion 
    after reaching the start.
\end{enumerate}
\textbf{Issues}:
\begin{itemize}
    \item We didn't run into any particular issues this week.
\end{itemize}
\textbf{Wave Code}:
For the wave code we created a new package called wave. The full details can be
seen in the \href{https://github.com/LeoBJenkins/GLaDOS}{repo}. Here is the code
for the wave demo node and trajectory:
\small
\verbatiminput{../../wave/wave/wave.py}
\end{document}